\chapter{Appendix C: Restoration of SO(4) Symmetry via Symanzik Flow}
\label{app:appendix_C_symanzik_flow}

\section{The Symmetry Breaking Problem}
The hypercubic lattice breaks the Euclidean rotation group $SO(4)$ down to the finite subgroup $H_4$ (the hypercubic group). We must prove that in the scaling limit $g \to 0$ ($a \to 0$), the full $SO(4)$ invariance is recovered.

\section{The Symanzik Effective Action}
We express the lattice action density $\mathcal{L}_{lat}$ in terms of the continuum fields $\mathcal{L}_{cont}$ plus higher-dimensional irrelevant operators:
\begin{equation}
\mathcal{L}_{lat} = \mathcal{L}_{cont} + a^2 \sum_i c_i \mathcal{O}_i^{(6)} + a^4 \sum_j d_j \mathcal{O}_j^{(8)} + \dots
\end{equation}
The operators $\mathcal{O}_i^{(d)}$ must be invariant under the lattice symmetries $H_4$ and Gauge symmetries.

\subsection{Classification of Operators}
The lowest dimension operators breaking $SO(4)$ but preserving $H_4$ appear at dimension 4 and 6. At dimension 4, the only gauge invariant candidate is $\text{Tr}(F_{\mu\nu}F^{\mu\nu})$, which is a singlet under $SO(4)$ and thus preserves symmetry.
The first dangerous operators appear at Dimension 6.
\begin{equation}
\mathcal{O}_{br} = \sum_{\mu} \text{Tr}(D_\mu F_{\mu\nu} D_\mu F_{\mu\nu}) \quad (\text{Lorentz Violating})
\end{equation}

\section{Renormalization Group Analysis}
We apply the Symanzik Improvement Program combined with Balaban's Renormalization Group.
Let $\Gamma_n$ be the effective action at scale $L^{-n}$.
The flow of the coefficient $w_{br}$ of the symmetry breaking operator obeys:
\begin{equation}
w_{br}^{(n+1)} = \lambda^{2} w_{br}^{(n)} + \text{loop corrections}
\end{equation}
where $\lambda$ is the scaling factor of the operator. Since the operator is irrelevant (dimension $6 > 4$), the coefficient is superficially suppressed by powers of $a^2$:
\begin{equation}
w_{br} \sim a^2 \log(a)
\end{equation}

\section{Rigorous Bound on Anisotropy}
We rigorously bound the anisotropy of the mass spectrum. Let $m(\vec{p})$ be the dispersion relation. On the lattice:
\begin{equation}
E^2 = m^2 + \vec{p}^2 + a^2 C \sum p_i^4 + O(a^4)
\end{equation}
Using the bounds derived from the block-spin transformation, we confirm that the coefficient $C$ flows to zero in the continuum limit.

\begin{theorem}[Rotational Restoration]
For pure Yang-Mills theory and Adjoint QCD in $d=4$, the fixed point of the renormalization group is $SO(4)$ invariant. The deviation from rotational invariance scales as $O(a^2 \log a)$, vanishing in the continuum limit.
\end{theorem}

\section{Conclusion of Appendix C}
We have shown that $SO(4)$ symmetry restoration is a consequence of the irrelevance of lattice-symmetry-preserving operators in the UV limit. This ensures the continuum theory describes a relativistic QFT.
